% Section: Positive definiteness of the Hessian on T_{2g}
%
% CONTEXT: This is the missing piece of the proof of the main theorem in the paper
% "D and P locally minimize Gauss variance." We need to show positive definiteness
% of the Hessian of J_8 on the irreducible representation T_{2g}.
%
% ============================================================
% SETUP (from the paper, fully self-contained)
% ============================================================
%
% Let p_1,...,p_8 be the vertices of a cube on S^2, in standard spherical
% coordinates (theta, phi):
%   p_1 = (theta_0, 0),       p_5 = (pi - theta_0, 0),
%   p_2 = (theta_0, pi/2),    p_6 = (pi - theta_0, pi/2),
%   p_3 = (theta_0, pi),      p_7 = (pi - theta_0, pi),
%   p_4 = (theta_0, 3*pi/2),  p_8 = (pi - theta_0, 3*pi/2),
% where theta_0 = arccos(1/sqrt(3)).
%
% Define the functional:
%   J_8 = I^+ * I^-,
% where
%   I^+ = integral_{S^2} exp(u) dA_{sph},
%   I^- = integral_{S^2} exp(-u) dA_{sph}  (improper integral, convergent),
% and
%   u = 2*pi * sum_{i=1}^{8} G_i(z),
%   G_i(z) = -(1/(2*pi)) * log(sin(dist(z, p_i)/2))
% is the Green function on S^2 satisfying
%   Delta_{FS} G_i = delta_{p_i} - 1/(4*pi).
%
% The symmetry group G of the cube acts on the 13-dimensional space
%   V = (direct sum_{i=1}^{8} T_{p_i} S^2) / SO(3).
% The Hessian of log J_8 at the cubic configuration decomposes under G into
% irreducible representations T_{1u} (+), E_g (+), E_u (-), T_{2u} (+), T_{2g} (?).
%
% For any deformation v = (v_1,...,v_8) with v_i in T_{p_i} S^2, at a critical point:
%
%   Hess_p log I^pm (v, v)
%   = integral_{S^2} (sum_i <2*pi nabla_z G_i, tilde{v}_i>)^2 d mu^pm
%     - integral_{S^2} sum_i <nabla_z u, tilde{v}_i> <2*pi nabla_z G_i, tilde{v}_i> d mu^pm,
%
% where mu^pm are the probability measures with densities exp(+/-u)/I^pm, and
% tilde{v}_i is any Killing field with tilde{v}_i(p_i) = v_i.
%
% For deformations of the "partitioned Killing field" type where each v_i equals
% either +partial_phi|_{p_i} or -partial_phi|_{p_i}, this simplifies to:
%
%   Hess_p log I^pm (v, v) = -16 pi^2 * integral_{S^2} (partial_phi G_A)(partial_phi G_B) d mu^pm,
%
% where A and B are the two index sets (with v_i = +partial_phi/sin(theta) for i in A,
% v_i = -partial_phi/sin(theta) for i in B), and G_A = sum_{i in A} G_i, G_B = sum_{i in B} G_i.
%
% Therefore:
%   Hess_p log J_8 (v, v) = -16 pi^2 * integral_{S^2} (partial_phi G_A)(partial_phi G_B) d(mu^+ + mu^-).
%
% The Hessian is positive definite on a deformation v iff this integral is NEGATIVE.
%
% ============================================================
% THE T_{2g} REPRESENTATION
% ============================================================
%
% The 3-dimensional irrep T_{2g} is generated (in the above coordinate system)
% by the orbit of the deformation:
%
%   v_i = +partial_phi/sin(theta)|_{p_i}  for i in A = {1, 2, 7, 8},
%   v_i = -partial_phi/sin(theta)|_{p_i}  for i in B = {3, 4, 5, 6}.
%
% Explicitly:
%   A = {p_1=(theta_0,0), p_2=(theta_0,pi/2), p_7=(pi-theta_0,pi), p_8=(pi-theta_0,3pi/2)},
%   B = {p_3=(theta_0,pi), p_4=(theta_0,3pi/2), p_5=(pi-theta_0,0), p_6=(pi-theta_0,pi/2)}.
%
% KEY OBSERVATION: B = R_pi(A), where R_pi: (theta,phi) |-> (theta, phi+pi)
% is the rotation by pi about the z-axis. Therefore:
%   G_B(theta,phi) = G_A(theta, phi+pi),
%   partial_phi G_B(theta,phi) = partial_phi G_A(theta, phi+pi).
%
% The measure mu^pm is invariant under R_pi (since the full cube is R_pi-symmetric).
%
% Hence the integral to be shown negative is:
%
%   I = integral_{S^2} (partial_phi G_A)(theta,phi) * (partial_phi G_A)(theta, phi+pi)
%         d(mu^+ + mu^-)(theta,phi).
%
% ============================================================
% PROBLEM TO PROVE
% ============================================================
%
% Prove that:
%
%   integral_{S^2} (partial_phi G_A)(theta,phi) * (partial_phi G_A)(theta, phi+pi)
%       d(mu^+ + mu^-)(theta,phi)  <  0,
%
% where:
%   - G_A = G_1 + G_2 + G_7 + G_8 is the sum of Green functions at
%     A = {(theta_0,0), (theta_0,pi/2), (pi-theta_0,pi), (pi-theta_0,3pi/2)},
%     with theta_0 = arccos(1/sqrt(3)).
%   - The probability measures mu^pm have densities exp(+/- u) / I^pm with
%     u = 2*pi*(G_1 + ... + G_8).
%   - partial_phi G_A is the derivative of G_A with respect to phi.
%
% Equivalently (using Fourier decomposition in phi): Writing
%   partial_phi G_A(theta, phi) = sum_m c_m(theta) e^{i m phi},
% the integral I equals
%   sum_m (-1)^m * integral_0^pi |c_m(theta)|^2 d nu(theta)
% for some positive measure nu (coming from the phi-integration of mu^+/- and
% the orthogonality). Since the cube has 4-fold symmetry, only m = 0, +-4, +-8,...
% and m = +-2, +-6,... and m = +-1, +-3, +-5, +-7,... modes can potentially survive,
% but the R_pi invariance of mu^pm means only modes with m even contribute.
% Actually: by the C_4 symmetry of the full configuration, only m = 0, +-4, +-8,...
% contribute to the average over phi. The (-1)^m factor is +1 for m divisible by 4
% and -1 for m = +-2, +-6, +-10, ... The question is whether the negative modes
% (m = 2 mod 4) dominate.
%
% HINTS FROM THE PAPER (for other irreps):
%
% For T_{2u} (different split A={2,3,4}, B={6,7,8} in different coordinates):
%   The integrand partial_phi G_A * partial_phi G_B was pointwise <= 0 on S^2,
%   giving positivity immediately.
%
% For T_{1u}, E_g (A/B are each a union of an upper ring half and lower ring half):
%   Used the Lemma about functions odd in phi with specific monotonicity to conclude
%   all cross-terms H_{ij}^pm = integral partial_phi G_i * partial_phi G_j d mu^pm
%   are negative.
%
% For T_{2g}: Direct sign argument fails because the 4 "vertical neighbor" pairs
%   (i,j) = (1,5), (2,6), (7,3), (8,4) give H_{ij}^pm > 0.  We need a combined
%   argument showing the total integral is still negative.
%
% SUGGESTED APPROACH (to try first):
%   1. Use the Fourier decomposition in phi. G_A(theta,phi) depends on the 4 points
%      {(theta_0,0),(theta_0,pi/2),(pi-theta_0,pi),(pi-theta_0,3pi/2)}.
%      By the symmetry phi -> phi + pi/2 sending A -> {(theta_0,pi/2),(theta_0,pi),
%      (pi-theta_0,3pi/2),(pi-theta_0,0)} (not A), G_A does not have 4-fold symmetry.
%      However, compute the Fourier modes explicitly and show the m=2 modes dominate.
%
%   2. Alternatively, use a comparison/monotonicity argument analogous to the Lemma
%      in the paper but adapted to the T_{2g} geometry.
%
%   3. Or: express the integral as a sum involving the full 8x8 matrix H_{ij}^pm
%      and use the explicit values (from the Lemma and direct computations) to
%      show the negative terms outweigh the positive ones.
%
% Please produce a complete, rigorous proof (in LaTeX) that the Hessian of log J_8
% is positive definite on T_{2g}, to be inserted into the paper.
\end{document}
